% Packages
% \usepackage[LGR,T1]{fontenc}
% \usepackage{pxfonts}
\usepackage[turkish,swedish,dutch,german,french,english,latin,irish,italian]{babel}    % Non aggiungere il giapponese, rompe la sillabazione in italiano.
\usepackage[expansion=all,protrusion=true]{microtype}
\usepackage{titlesec}
\usepackage{geometry}
\usepackage{fancyhdr}
\usepackage[lines=3, loversize=0.12, lraise=0]{lettrine}
\usepackage{hyperref}
\usepackage{xspace}
\usepackage{fontspec}
\usepackage{soul}                   % per l'evidenziatore
\usepackage{xcolor}                 % per i colori
\usepackage{enumitem}
\usepackage{graphicx}               % per le illustrazioni
\usepackage{adjustbox}
\usepackage[indentfirst=false]{quoting}
\usepackage{afterpage}
\usepackage{tocloft}
\usepackage{hyphenat}
\usepackage{ifthen}
\usepackage{pdftexcmds}             % per la condizionale sulla proprietà \purpose
\usepackage{ragged2e}

% Dimensioni pagina
\geometry{%
    paperwidth=15cm,%
    paperheight=21cm,%
    top=2.8cm,%
    bottom=3.3cm,%
    inner=2cm,%
    outer=1.8cm,%
}

\raggedbottom

% Capitoli: solo numero romano + titolo sulla stessa riga
\renewcommand{\thechapter}{\Roman{chapter}}
\setlength{\cftchapnumwidth}{3.6em}

\titleformat{\chapter}% Il comando di sezionamento
{\normalfont\Large\scshape\centering}% Stile per l'intero titolo
{}%
{0pt}% Spazio tra numerazione e testo
{\thechapter. }% Prefisso al testo del titolo

\titlespacing*{\chapter}{0pt}{50pt}{40pt}

% HEADER / FOOTER
\pagestyle{fancy}
\fancyhf{}                                      % pulisce tutto
\renewcommand{\headrulewidth}{0pt}              % nessuna linea sotto l'intestazione
\fancyhead[RO]{\textsc{\leftmark}}              % pagina dispari: capitolo in maiuscoletto
\fancyhead[LE]{\textsc{\mytitle}}               % pagina pari:    titolo libro in maiuscoletto

\fancypagestyle{plain}{%
    \fancyhf{}                                  % pulisce tutto
    \renewcommand{\headrulewidth}{0pt}          % nessuna linea
    \fancyfoot[LE,RO]{\thepage}                 % numero pagina all’esterno ANCHE sulla prima pagina del capitolo
}

\fancyfoot[LE,RO]{\thepage}                     % Footer numero pagina all'esterno

% Elimina "Chapter 3", "Capitolo II", corsivo e maiuscolo
\makeatletter
\renewcommand\chaptermark[1]{\markboth{#1}{#1}} % solo il titolo del capitolo
\renewcommand\sectionmark[1]{\markright{#1}}    % (se usi sezioni)
\makeatother

% Usiamo \rightmark per il titolo del libro (lo impostiamo una volta)
\markright{\mytitle}

\selectlanguage{italian}

% Rimuoviamo tutti i bordi colorati e mettiamo una sottolineatura punteggiata
\hypersetup{%
    colorlinks  = false,                        % false = usa i box (ma noi li ridefiniamo)
    pdfborder  = {0 0 0},                       % elimina completamente il bordo standard
% Stile personalizzato per tutti i link interni
    linkbordercolor = {white},                  % inutile, ma per sicurezza
}

\xspaceaddexceptions{\dots}
\xspaceaddexceptions{»}
\xspaceaddexceptions{]}

% Definiamo il nuovo stile dei link: sottolineatura punteggiata (dotted)
% Invece di usare colorlinks=true (che mette il colore al testo),
% usiamo pdfborderstyle per ottenere una sottolineatura punteggiata

\hypersetup{%
    pdfborderstyle = {/S/D /W 1 /D [3 2]}       % S = underlined, D = dashed, W = spessore 1pt, pattern [3 2]
}

% Se preferisci la sottolineatura punteggiata ma in nero (più sobria):
% \hypersetup{pdfborderstyle={/S/D /W 1 /D [3 2]}, linkbordercolor={black}}

% --------------------- NOTE A PIÈ DI PAGINA IN CORSIVO (stesso corpo) ---------------------
\makeatletter
\renewcommand{\@makefntext}[1]{%
    \parindent 1em\noindent
    \hb@xt@1.8em{\hss\@makefnmark}%
    \normalfont\itshape #1%
}
\makeatother

\setlist{%
    parsep   = 0pt,                             % spazio tra paragrafi dentro un item (se usi \par)
    itemsep  = 0pt,                             % spazio verticale tra un item e l’altro ← QUESTO È IL PRINCIPALE
}

% \setmainfont{BaskervilleF}[%                    % Richiesto per small caps
%     Extension       = .otf,%
%     UprightFont     = *-Regular,%
%     ItalicFont      = *-Italic,%
%     BoldFont        = *-Bold,%
%     BoldItalicFont  = *-BoldItalic,%
% ]
% \setlength{\headheight}{14pt}                   % Evita un warning per ogni pagina

% \setmainfont{EBGaramond}
% \microtypesetup{tracking={true, smallcaps}}   % For Garamond only
\setmainfont{EBGaramond}[
    Path            = fonts/EB_Garamond/static/,
    Extension       = .ttf,
    UprightFont     = *-Regular,
    ItalicFont      = *-Italic,
    BoldFont        = *-Medium,
    BoldItalicFont  = *-MediumItalic,
]

\setlength{\parskip}{0pt}
\setlength{\parindent}{0.3cm}

\renewcommand{\cfttoctitlefont}{\relax}         % non stampa nulla
\setlength{\cftbeforetoctitleskip}{-4em}        % elimina spazio sopra (regola se serve)
\setlength{\cftaftertoctitleskip}{0em}          % elimina spazio sotto
\renewcommand{\contentsname}{}                  % titolo vuoto
\renewcommand{\cftchapfont}{\normalfont}        % capitoli
\renewcommand{\cftchappagefont}{\normalfont}    % capitoli

% Comandi personalizzati
\newcommand{\threestars}{%
    \par
    \addvspace{1.2\baselineskip}%
    {\centering\Large\bfseries***\par}
    \addvspace{0.8\baselineskip}%
%    \noindent%
}

% Separatore per le sezioni di entrelacement
\newcommand{\entrelacement}{\threestars}

% Separatore per gli stacchi temporali
\newcommand{\temporaljump}{\bigskip}
% \newcommand{\temporaljump}{\bigskip\noindent}

\def\currentlanguage{italian}

\makeatletter
\newcommand{\myforeignlanguage@format@book}[2]{\foreignlanguage{#1}{#2}}%
\newcommand{\myforeignlanguage@format@spellcheck}[3]{ @#1 #2 @#3}% Mantieni uno spazio prima di #1, per spezzare l'eventuale apostrofo precedente.
\newcommand{\myforeignlanguage}[2]{%
    \ifnum\pdf@strcmp{\detokenize\expandafter{\purpose}}{spellcheck}=0%
        \protect\myforeignlanguage@spellcheck@safe{#1}{#2}%
    \else
        \myforeignlanguage@format@book{#1}{#2}%
    \fi
}%
\newcommand{\myforeignlanguage@spellcheck@safe}[2]{%
    \begingroup%
    \edef\previouslanguage{\currentlanguage}%
    \xdef\currentlanguage{#1}%
    \myforeignlanguage@format@spellcheck{#1}{#2}{\previouslanguage}%
    \xdef\currentlanguage{\previouslanguage}%
    \endgroup%
}%
\makeatother

\newcommand{\dutchname}[1]{\myforeignlanguage{dutch}{#1}}%
\newcommand{\dutch}[1]{\textit{\myforeignlanguage{dutch}{#1}}}%
\newcommand{\englishname}[1]{\myforeignlanguage{english}{#1}}%
\newcommand{\english}[1]{\textit{\myforeignlanguage{english}{#1}}}%
\newcommand{\frenchname}[1]{\myforeignlanguage{french}{#1}}%
\newcommand{\french}[1]{\textit{\myforeignlanguage{french}{#1}}}%
\newcommand{\germanname}[1]{\myforeignlanguage{german}{#1}}%
\newcommand{\german}[1]{\textit{\myforeignlanguage{german}{#1}}}%
\newcommand{\irishname}[1]{\myforeignlanguage{irish}{#1}}%
\newcommand{\irish}[1]{\textit{\myforeignlanguage{irish}{#1}}}%
\newcommand{\japanesename}[1]{\mbox{#1}}%
\newcommand{\japanese}[1]{\textit{\mbox{#1}}}%
\newcommand{\latinname}[1]{\myforeignlanguage{latin}{#1}}%
\newcommand{\latin}[1]{\textit{\myforeignlanguage{latin}{#1}}}%
\newcommand{\provencalname}[1]{\myforeignlanguage{french}{#1}}%
\newcommand{\provencal}[1]{\textit{\myforeignlanguage{french}{#1}}}% FIXME
\newcommand{\swedishname}[1]{\myforeignlanguage{swedish}{#1}}%
\newcommand{\swedish}[1]{\textit{\myforeignlanguage{swedish}{#1}}}%
\newcommand{\turkishname}[1]{\myforeignlanguage{turkish}{#1}}%
\newcommand{\turkish}[1]{\textit{\myforeignlanguage{turkish}{#1}}}%

\newcommand{\yellow}[1]{{\sethlcolor{yellow!25}\hl{#1}}}%

\newcommand{\blankpage}{%
    \clearpage
    \thispagestyle{empty}%
    \null
    \clearpage
}%

\newboolean{illustrations}
\setboolean{illustrations}{true}

\newcommand{\illustration}[1]{%
    \ifthenelse{\boolean{illustrations}}{%
        \begin{center}%
        \vspace*{\fill} % spazio elastico sopra
        \adjustbox{max size={1\textwidth}{1\textheight}, center}{\includegraphics{#1}}%
        \vspace*{\fill} % spazio elastico sotto (uguale a quello sopra)
        \end{center}%
        \clearpage
    }{}
}%

\newcommand{\fullpageillustration}[2][1\textwidth]{%
    \ifthenelse{\boolean{illustrations}}{%
        \afterpage{%
            \clearpage
            \thispagestyle{empty}% or plain if you prefer page number
            \null
            \vfill
            \begin{center}%
            \adjustbox{max size={#1}{0.86\textheight},center}{\includegraphics[width=#1]{#2}}%
            \end{center}%
            \vfill
            \clearpage
        }%
    }{}
}%

\newcommand{\afterpageillustration}[2][0.95\textwidth]{%
    \ifthenelse{\boolean{illustrations}}{%
        \afterpage{%
            \clearpage
            \thispagestyle{empty}%
            \null
            \vfill
            \centering
            \adjustbox{max size={#1}{#1}, center}{%
                \includegraphics{#2}%
            }%
            \vfill
            \clearpage
        }%
    }{}
}%

\newcommand{\afterpagetopillustration}[2][0.95\textwidth]{%
    \ifthenelse{\boolean{illustrations}}{%
        \afterpage{%
            \clearpage
            \noindent
            \centering
            \adjustbox{max size={#1}{\textheight}, center}{%
                \includegraphics{#2}%
            }%
            \vspace{3em}%
            \par
        }%
    }{}
}%

\newcommand{\poetry}[1]{%
    \begin{quoting}%
    \itshape
    \begingroup
    #1%
    \endgroup
    \end{quoting}%
}%

\newcommand{\quotationx}[2][\itshape]{%
    \begin{quoting}%
    #1%
    \begingroup
    #2%
    \endgroup
    \end{quoting}%
}%

\usepackage{relsize}

\newcommand{\newspaper}[1]{%
    \begin{quoting}%
    \itshape%
    \smaller[0.5]%
    \begingroup
    #1%
    \endgroup
    \end{quoting}%
}%

\makeatletter
\newcommand{\email@format@audio}[5]{%
    \begingroup%
    \ifx\relax#5\relax%
        #3. {#4}%
    \else%
        \spoken{#5}{#3. {#4}}{}%
    \fi%
    \endgroup%
}%
\newcommand{\email@format@book}[4]{%
    \begin{quoting}%
    \begingroup
    \smaller[0.5]%
    \englishname{To}: #2\\
    \englishname{From}: #1\\
    \englishname{Subject}: #3\\
    {#4}%
    \endgroup
    \end{quoting}%
}%

\newcommand{\email}[5]{%
    \ifnum
        \pdf@strcmp{\detokenize\expandafter{\purpose}}{audio}=0 %
        \email@format@audio{#1}{#2}{#3}{#4}{#5}
    \else
        \email@format@book{#1}{#2}{#3}{#4}
    \fi
}

\newcommand{\note@format@audio}[1]{}%
\newcommand{\note@format@book}[1]{%
    \begin{quoting}[leftmargin=0pt,rightmargin=0pt]%
    \color{red}%
    \small
    \itshape
    \begingroup
    #1%
    \endgroup
    \end{quoting}%
}%

\newcommand{\note}[1]{%
    \ifnum%
        \pdf@strcmp{\detokenize\expandafter{\purpose}}{audio}=0 %
        \note@format@audio{#1}%
    \else%
        \note@format@book{#1}%
    \fi%
}%
\makeatother

% marginpar
\newcommand{\comment}[1]{}%

\newcommand{\ellipsis}{\ldots\xspace}%  % FIXME: mette spazio anche prima del » chiuso

% Interruzione improvvisa di un discorso.
\newcommand{\interruption}{---}

\newcommand{\cover}[3]{%
    \thispagestyle{empty} % no header, footer, or page number
    \vspace*{\fill}%
    \begin{center}%
    {\normalfont\bfseries\centering\fontsize{27}{27}\selectfont #1}%
    \end{center}%
    \bigbreak
    \begin{center}%
    {#2}%
    \end{center}%
%    \vspace{1cm}%
    \begin{center}%
    \illustration{#3}%
    \end{center}%
    \vspace*{\fill}%
    \clearpage
    \thispagestyle{empty} % no header, footer, or page number
}

\newcommand{\disclaimer}{
    \thispagestyle{empty} % no header, footer, or page number
    \vspace*{\fill}
    \begin{flushleft}%
    \textit{%
        Questa è un'opera di fantasia.
        Nomi, personaggi, luoghi ed eventi sono prodotto dell'immaginazione dell’autore.
        Qualsiasi somiglianza con persone reali, vive o morte, con eventi o luoghi esistenti è puramente casuale.
    }
    \end{flushleft}%
    \clearpage
}

\newcommand{\epigraph}[1]{%
    \cleardoublepage
    \thispagestyle{empty} % no header, footer, or page number
    \vspace*{\fill}%
    \begin{flushright}%
    \textit{#1}%
    \end{flushright}%
    \vspace*{\fill}%
    \clearpage
    \thispagestyle{empty} % no header, footer, or page number
}

% Comandi per i virgolettati
% Variabile "globale" (ma meglio limitata al gruppo quando possibile)
\newtoks\spokenAnteToks   % usiamo toks per evitare problemi con \edef/global

% Resetta sempre all'inizio di un paragrafo o come preferisci
%%%% \AtBeginDocument{\spokenAnteToks={}}   % o in \everypar, se serve

\global\spokenAnteToks={}%

\usepackage{xstring}

\newcounter{dropcaplines}
\setcounter{dropcaplines}{3}

\makeatletter
\newcommand{\dropcap@format@pdf}[2]{\lettrine[lines=\value{dropcaplines}]{#1}{#2}}%
\newcommand{\dropcap@format@other}[2]{#1#2}%
\newcommand{\dropcap}[2]{
    \ifnum\pdf@strcmp{\detokenize\expandafter{\purpose}}{pdf}=0 %
        \dropcap@format@pdf{#1}{#2}%
    \else
        \dropcap@format@other{#1}{#2}%
    \fi
}%
\makeatother

\makeatletter
\newcommand{\mylettrine@format@pdf}[3]{\lettrine[lines=\value{dropcaplines},ante=#3]{#1}{#2}}%
\newcommand{\mylettrine@format@other}[2]{#1#2}%
\newcommand{\mylettrine}[3]{
    \ifnum\pdf@strcmp{\detokenize\expandafter{\purpose}}{pdf}=0 %
        \mylettrine@format@pdf{#1}{#2}{#3}%
    \else
        \mylettrine@format@other{#1}{#2}%
    \fi
}%
\makeatother

\NewDocumentCommand{\smartdropcap}{m}{%
    \begingroup
    \let\xspace\empty
    \edef\temp{#1}%
    \expandafter\endgroup
    \expandafter\smartdropcapaux\expandafter{\temp}%
}

\NewDocumentCommand{\smartdropcapaux}{m}{%
    \edef\arg{#1}%
    \edef\firstletter{\directlua{tex.sprint(unicode.utf8.sub("\arg",1,1))}}%
    \edef\remainder{\directlua{tex.sprint(unicode.utf8.sub("\arg",2))}}%
    \mylettrine{\firstletter}{\MakeLowercase{\remainder}}{\the\spokenAnteToks}%
}

\makeatletter
\newcommand{\acronym@format@book}[1]{\textsc{\lowercase{{#1}}}}%
\newcommand{\acronym}[1]{%
    \ifnum\pdf@strcmp{\detokenize\expandafter{\purpose}}{pdf}=0 %
        \acronym@format@book{#1}%
    \else%
        {#1}%
    \fi%
}%

\newcommand{\spoken@format@book}[1]{%
    \begingroup
    \edef\temparg{\detokenize{#1}}%
    \IfBeginWith{\temparg}{\detokenize{\smartdropcap}}{\global\spokenAnteToks={«}}{«}%
    \expandafter\endgroup
    #1»%
    \global\spokenAnteToks={}%
}
\newcommand{\spoken@format@audio}[3]{\par[#1\ifx\relax#3\relax\else, #3\fi]: #2\par}%
\newcommand{\spoken}[3]{%
    \ifnum\pdf@strcmp{\detokenize\expandafter{\purpose}}{audio}=0 %
        \spoken@format@audio{#1}{#2}{#3}%
    \else%
        \spoken@format@book{#2}%
    \fi%
}%

\newcommand{\thought@format@book}[1]{``#1''}%
\newcommand{\thought@format@audio}[3]{\par[#1\ifx\relax#3\relax\else, #3\fi]: #2\par}%
\newcommand{\thought}[3]{%
    \ifnum\pdf@strcmp{\detokenize\expandafter{\purpose}}{audio}=0 %
        \thought@format@audio{#1}{#2}{#3}%
    \else%
        \thought@format@book{#2}%
    \fi%
}
\makeatother

% Comando base con parametri utili
\newcommand{\letterstyle}{%
    \noindent
    \RaggedRight
    \leftskip=1cm          % margine sinistro (modifica a piacere)
    \parindent=0pt
    \hyphenpenalty=10000     % disabilita la sillabazione (valore molto alto)
    \exhyphenpenalty=10000   % anche per sillabazione forzata
    \pretolerance=9999
    \tolerance=9999
    \emergencystretch=3em
%    \setstretch{2}
}

\newcommand{\letter}[1]{%
    \begingroup%
    \bigskip
    \letterstyle
    \itshape
    \par#1\par%
    \bigskip
    \endgroup%
}

\newcommand{\newspapername}[1]{#1}%

\newcommand{\sxx}[2][]{\spoken{default}{#2}{#1}}%
\newcommand{\sxa}[2][]{\spoken{senza nome a}{#2}{#1}}%
\newcommand{\sxb}[2][]{\spoken{senza nome b}{#2}{#1}}%
\newcommand{\sxc}[2][]{\spoken{senza nome c}{#2}{#1}}%
\newcommand{\sxd}[2][]{\spoken{senza nome d}{#2}{#1}}%
\newcommand{\sxe}[2][]{\spoken{senza nome e}{#2}{#1}}%
\newcommand{\sxf}[2][]{\spoken{senza nome f}{#2}{#1}}%
\newcommand{\txx}[2][]{\thought{default}{#2}{#1}}%
% citazione diretta breve
\newcommand{\sq}[1]{``#1''}%

\newcommand{\Pffft}{Pffft\ellipsis}
\newcommand{\Mmmh}{Mmmh\ellipsis}
\newcommand{\beh}{be'}
\newcommand{\Beh}{Be'}
